\section{Value-Specific Identification}
\label{sec:identifiability}

Section~\ref{sec:formulation} showed how to calculate protocol value when the
latent covariance $K$ is known. We now ask what can be learned when the only
available population information is the joint law of the measurements and
target under the realised protocol $A$. Throughout this section, $A$ and $B$
denote the measurement operators induced by the realised and alternative
protocols.
Sections~\ref{sec:estimation}--\ref{sec:design} return to the action-set notation
$S$, with measurement matrix $A_S$.

We begin with the linear target $\Theta=h^\top Z$. Then $(Y_A,\Theta)$ is jointly
Gaussian, so its full law is determined by second moments. We abbreviate
$\cI(B;K)=\cI_g(B;K)$ for
$g(z)=z$; in this linear-Gaussian case it also equals $\cI_{\mathrm L}(B)$.
We index grid points by $0,\dots,p-1$. Nonlinear aggregates require more than
moment matching; we treat them separately through an exact permutation argument
in \cref{prop:permutation}.

The law of $(Y_A,\Theta)$ may fail to identify $K$ even when it determines the
scalar value of the specified alternative $B$.

\subsection{Identification Criterion}
\label{sec:value-specific-identification}

For a benchmark law $\Prob_A$, define its observational equivalence class by
\begin{equation}
\mathcal K_A(\Prob_A)=\{K:\Prob_K(Y_A,\Theta)=\Prob_A\}.
\label{eq:equivalence-class}
\end{equation}
Every $K\in\mathcal K_A(\Prob_A)$ generates exactly the same distribution of
the benchmark data. No amount of additional sampling under the same protocol
$A$ can distinguish them. Consequently, for any target and model class under
consideration,
\begin{equation}
\boxed{\begin{gathered}
\cI_g(B;\cdot)\ \text{is identified from }\Prob_A\\[-1pt]
\Longleftrightarrow\quad
\cI_g(B;K_1)=\cI_g(B;K_2)
\quad\text{for all }K_1,K_2\in\mathcal K_A(\Prob_A).
\end{gathered}}
\label{eq:value-specific-criterion}
\end{equation}
This criterion is value-specific: it requires uniqueness only of the specified
alternative's value, not of every component of $K$. The class
$\mathcal K_A(\Prob_A)$ may contain many covariance structures while assigning
a unique value to $B$. Conversely, two members of this class with different
values establish non-identification. This is functional identification over an
observational-equivalence class in the sense of
\citet{koopmans1950identification}.

For the linear target $\Theta=h^\top Z$, the realised benchmark law reveals
three covariance blocks. Take $h=\omega$, and let
$Y_A=AZ+\varepsilon_A$ with $\varepsilon_A\sim\cN(0,R_A)$ and $R_A$ known.
Then $(Y_A,\Theta)$ is jointly Gaussian with covariance
\begin{equation}
\begin{pmatrix}
AKA^\top+R_A & AKh\\
h^\top KA^\top & h^\top Kh
\end{pmatrix}.
\label{eq:observables}
\end{equation}
The Gaussian benchmark law therefore depends on $K$ through the three blocks
\begin{equation}
AKA^\top,\qquad AKh,\qquad h^\top Kh,
\label{eq:three-blocks}
\end{equation}
which constitute its complete covariance information about $K$. The first
block is the covariance of the realised measurements after subtracting the
known measurement noise. The second is their cross-covariance with the target,
and the third is the target variance. Any change in $K$ that leaves all three
blocks unchanged is invisible to the benchmark.



\subsection{Minimal Stationary Counterexample}

We first ask when invisible dependence can arise in the simplest stationary
setting. A one-point benchmark reveals no lag information through
$\Var(Y_A)$, because the process is standardised. It reveals only two weighted
sums of the lag correlations: one through
$\Cov(Y_A,\Theta)$ and one through $\Var(\Theta)$. These two equations determine
all unknown lags on grids of at most three points. On a four-point grid they
leave one degree of freedom for the first time, and moving along that direction
can change the value of a protocol that observes $(Z_1,Z_2)$.

\begin{theorem}[Sharp minimal stationary counterexample]
\label{thm:minimal}
Let $p\ge2$ and let $Z$ be stationary and standardised on the grid
$\{0,1,\dots,p-1\}$ with
correlations $\rho(1),\dots,\rho(p-1)$, let $\Theta=p^{-1}\sum_j Z_j$, and let
$A$ observe $Z_0$ with known noise variance $\nu^2$. Then
\begin{enumerate}[label=(\roman*),leftmargin=2.2em,itemsep=2pt]
\item for $p\le3$ the map from $(\rho(1),\dots,\rho(p-1))$ to the observable
functionals is injective, so the stationary lag-correlation vector, and hence
$K$ within this model class, is identified;
\item for $p\ge4$ the space of stationary invisible directions has dimension
$p-3\ge1$.
\end{enumerate}
For $p=4$, the kernel is spanned by the lag perturbation
\begin{equation}
\delta=(\delta(0),\delta(1),\delta(2),\delta(3))=(0,1,-2,1).
\label{eq:delta4}
\end{equation}
For any positive-definite base correlation matrix $K$ with lag vector $\rho$,
let $K_\pm$ have lag vectors $\rho\pm\varepsilon\delta$, with lag $0$ fixed at
one. For all sufficiently small $\varepsilon>0$, these matrices are admissible
and observationally equivalent under $A$. If $B$ observes $(Z_1,Z_2)$ with
independent measurement noise $R_B=\nu_B^2 I_2$, where $\nu_B^2\ge0$, then
\begin{equation}
\cI(B;K_\pm)=\frac{2b^2}{\{1+\nu_B^2+\rho(1)\pm\varepsilon\}\,\Var(\Theta)},
\qquad
b=\frac{1+2\rho(1)+\rho(2)}{4},
\label{eq:gap-closed-form}
\end{equation}
so the two values differ whenever $1+2\rho(1)+\rho(2)\neq0$.
\end{theorem}

\begin{proof}[Proof idea]
The realised point and the target reveal two independent weighted sums of the
$p-1$ unknown lag correlations. Their Jacobian has rank
$\min\{2,p-1\}$, leaving $p-3$ free directions exactly when $p\ge4$; at
$p=4$, solving the two equations gives $(1,-2,1)$. Along this direction the
target variance and its covariance with $Z_1$ and $Z_2$ remain fixed, whereas
the covariance of $(Z_1,Z_2)$ changes, which yields
\eqref{eq:gap-closed-form}. The calculation is given in
Appendix~\ref{sec:app-identification}.
\end{proof}

This four-point threshold is specific to the stationary, standardised,
one-observation model with a uniform mean target. \Cref{fig:identifiability}
uses $p=4$, $h=\tfrac{1}{4}\mathbf{1}_4$, the baseline profile
$\rho_0(u)=\exp(-u)$ (that is, $\tau=1$), and the raw lag perturbation
$\varepsilon=\numIdEps$ in \eqref{eq:delta4}. Both protocols are noiseless,
with $A=e_0^\top$ and $B=(e_1,e_2)^\top$. The two Toeplitz correlation
matrices $K_\pm$ preserve $\Var(Y_A)=1$,
$\Cov(Y_A,\Theta)=\numIdCovYTheta$ and
$\Var(\Theta)=\numIdVarTheta$ to a maximum discrepancy of
\numIdDiscrepancy. Substitution into \eqref{eq:gap-closed-form} gives
$\cI(B;K_+)=\numIdCeilPlus$ and $\cI(B;K_-)=\numIdCeilMinus$.

\begin{figure}[t]
\centering
\includegraphics[width=0.72\linewidth]{figures/fig_identifiability.pdf}
\caption{Four-point instance of \cref{thm:minimal}. The profiles
$\rho_\pm=\rho_0\pm\varepsilon\delta$ are observationally equivalent under
$A=\{Z_0\}$ but assign values \numIdCeilPlus{} and \numIdCeilMinus{} to
$B=\{Z_1,Z_2\}$.}
\label{fig:identifiability}
\end{figure}

\FloatBarrier

\subsection{Invisible Directions}

The vector $(0,1,-2,1)$ is one instance of a more general object: a covariance
perturbation that changes latent dependence while preserving every covariance
block visible under $A$.

\begin{definition}[Invisible direction]
\label{def:invisible}
A symmetric $\Delta\in\R^{p\times p}$ with $\diag(\Delta)=0$ is
\emph{invisible} to $(A,h)$ if
\begin{equation}
A\Delta A^\top=0,\qquad A\Delta h=0,\qquad h^\top\Delta h=0 .
\label{eq:invisible}
\end{equation}
The three equalities preserve, respectively, the realised-measurement
covariance, the realised-measurement--target cross-covariance and the target
variance. The zero-diagonal requirement preserves the unit diagonal. For
positive-definite $K$, sufficiently small perturbations $K\pm\varepsilon\Delta$ remain
correlation matrices. We write $\cD(A,h)$ for the
linear space of directions invisible to $(A,h)$.
\end{definition}

\begin{theorem}[Value-changing invisible directions imply non-identification]
\label{thm:impossibility}
Let $0\neq\Delta\in\cD(A,h)$ and let $K_0\succ0$ be a correlation matrix. For
all sufficiently small $|\varepsilon|$, $K_\varepsilon=K_0+\varepsilon\Delta$
is a correlation matrix and has the same benchmark law as $K_0$; in
particular,
\begin{equation}
\Prob_{K_\varepsilon}(Y_A,\Theta)=\Prob_{K_0}(Y_A,\Theta).
\label{eq:obs-equivalent}
\end{equation}
Let $B$ be an alternative protocol with noise $R_B\succeq0$ such that
$\Sigma_B(K_0)=BK_0B^\top+R_B\succ0$, and set
$W=\Sigma_B(K_0)^{-1}$. Its directional derivative is
\begin{equation}
\begin{aligned}
D_B(\Delta;K_0)
&:=\left.\frac{\dd}{\dd\varepsilon}
\cI(B;K_0+\varepsilon\Delta)\right|_{\varepsilon=0}\\
&=\frac{2h^\top\Delta B^\top WBK_0h
-h^\top K_0B^\top W(B\Delta B^\top)WBK_0h}{h^\top K_0h}.
\end{aligned}
\label{eq:ceiling-derivative}
\end{equation}
If $D_B(\Delta;K_0)\neq0$, then the value functional $\cI(B;\cdot)$ is not
locally identified at $K_0$ within the model class: for every sufficiently
small $\varepsilon>0$, the observationally equivalent matrices
$K_0\pm\varepsilon\Delta$ assign different values to $B$.
\end{theorem}

\begin{proof}[Proof idea]
The three invisibility equations preserve every covariance block of the
Gaussian benchmark law. Differentiating the alternative-protocol value with
the inverse rule gives \eqref{eq:ceiling-derivative}; a non-zero derivative
then separates the two compatible paths $K_0\pm\varepsilon\Delta$ to first
order. The full matrix calculation appears in
Appendix~\ref{sec:app-identification}.
\end{proof}

Thus no estimator based only on $(Y_A,\Theta)$ can be consistent at both local
alternatives, and collecting more units under the same protocol does not
resolve the ambiguity.

If $A$ has $d$ rows, rank--nullity gives a lower bound on the dimension of the
invisible space. A standardised $K$ has $p(p-1)/2$ free entries whereas
\eqref{eq:three-blocks} supplies at most $d(d+1)/2+d+1$ linear constraints.
Rank--nullity therefore gives
\begin{equation}
\dim\cD(A,h)\ge
\frac{p(p-1)}2-\frac{d(d+1)}2-d-1,
\qquad
\text{when the right-hand side is positive.}
\label{eq:counting}
\end{equation}
Near a positive-definite $K_0$, the equivalence class contains the corresponding
open piece of the affine space $K_0+\cD(A,h)$
\citep{rothenberg1971identification}. This count only measures ambiguity in
$K$: value non-identification still requires sensitivity in
\eqref{eq:ceiling-derivative}. With fixed $d$, the contrast is nevertheless
stark---the benchmark contributes $O(d^2)$ constraints to $O(p^2)$ covariance
parameters.

\subsection{Nonlinear Targets}

The invisible-direction theorem uses joint Gaussianity of
$(Y_A,\Theta)$ and therefore applies directly to a linear target. With a
nonlinear aggregate, equality of the first two moments does not imply equality
of the full benchmark law. The following symmetry argument instead constructs
two exactly equivalent benchmark laws. The condition $A\Xi=A$ says that the
permutation leaves every realised measurement unchanged, while
$\Xi^\top\omega=\omega$ says that it preserves the aggregate weights.

\begin{proposition}[Permutation equivalence for arbitrary targets]
\label{prop:permutation}
Let $\Xi$ be a permutation matrix with $A\Xi=A$ and $\Xi^\top\omega=\omega$, and put
$K'=\Xi K\Xi^\top$. For every $g\in L^2(\phi)$ with $V_g(K)>0$, the joint laws of
$(Y_A,\Theta_g)$ under $K$ and under $K'$ coincide, and
\begin{equation}
\cI_g(B;K')=\cI_g(B\Xi;K)
\qquad\text{for every alternative protocol }B .
\label{eq:perm-transport}
\end{equation}
Consequently, whenever $\cI_g(B\Xi;K)\neq\cI_g(B;K)$, the value of $B$ is not
identified from the benchmark law.
\end{proposition}

Here $B\Xi$ is the algebraically permuted measurement operator used to compare
values, with the same observation-noise covariance $R_B$ as $B$; it need not
belong to the original candidate family.

\begin{proof}
Let $Z\sim\cN(0,K)$ and $Z'=\Xi Z\sim\cN(0,K')$. Then $AZ'=A\Xi Z=AZ$, and since
$g$ acts entrywise $g(\Xi Z)=\Xi g(Z)$, so
$\Theta_g(Z')=\omega^\top\Xi g(Z)=(\Xi^\top\omega)^\top g(Z)=\Theta_g(Z)$; in
particular $V_g(K')=V_g(K)>0$. The measurement noise is independent with the
same law under both, so the joint laws agree. Applying the same substitution to
$B$ gives \eqref{eq:perm-transport}.
\end{proof}

Non-identification follows whenever the permutation changes the value of $B$.
For example, no change occurs when $\Xi K\Xi^\top=K$. The next construction
shows that three grid points suffice for every nonconstant target under uniform
aggregation.

\begin{example}[Three points suffice under uniform aggregation]
\label{ex:three-point}
Take $p=3$, $A=e_0^\top$, uniform $\omega$, let $\Xi$ transpose coordinates $1$
and $2$, and let $B=e_1^\top$ observe $Z_1$ with noise $\nu_B^2$, writing
$s=1+\nu_B^2$. For $0<a<1$ put
\begin{equation}
K_a=\begin{pmatrix}1&a&0\\a&1&0\\0&0&1\end{pmatrix},
\qquad K_a'=\Xi K_a\Xi^\top .
\label{eq:three-point}
\end{equation}
Both are correlation matrices, $A\Xi=A$ and $\Xi^\top\omega=\omega$, so the
benchmark laws coincide. Nevertheless,
$\cI_g(B;K_a)>\cI_g(B;K_a')$ for every nonconstant
$g\in L^2(\phi)$, every $a\in(0,1)$ and every noise level. The explicit
calculation, which reduces the gap to positive values of $C_g$, is given in
Appendix~\ref{sec:app-identification}.
\end{example}

Without stationarity, three points suffice for arbitrary nonconstant
square-integrable transformations under uniform aggregation; within the
stationary standardised model with a one-point benchmark and a mean target, the
sharp minimum is four.

\subsection{Identification by Augmentation}
\label{sec:augmentation}

Suppose $A$ has already been run and a further acquisition $A'$ is contemplated
on the same units. Let $A_{\rm aug}$ stack the rows of $A$ and $A'$. Their joint
measurement reveals the covariance within $A'$, its cross-covariance with $A$,
and the corresponding target cross-covariances. Collect the revealed linear
directions in
\begin{equation}
\bar A=\begin{bmatrix}A_{\rm aug}\\ h^\top\end{bmatrix} .
\label{eq:abar}
\end{equation}
Let
\begin{equation}
\cD(A_{\rm aug},h)=\big\{\Delta=\Delta^\top:\diag(\Delta)=0,\;
A_{\rm aug}\Delta A_{\rm aug}^\top=0,\;
A_{\rm aug}\Delta h=0,\;h^\top\Delta h=0\big\}
\label{eq:joint-invisible}
\end{equation}
be the invisible space of the augmented benchmark.

\begin{proposition}[Exact sufficiency and a first-order failure certificate]
\label{prop:augmentation}
Fix $K_0\succ0$ with $\diag(K_0)=\Ind$ and an alternative protocol $B$ with
known $R_B\succeq0$ such that $BK_0B^\top+R_B\succ0$.
\begin{enumerate}[label=(\roman*),leftmargin=2.2em,itemsep=2pt]
\item \emph{Exact sufficiency.} If $\rowspace(B)\subseteq\rowspace(\bar A)$,
then the joint Gaussian law of $(Y_B,\Theta)$, and hence $\cI(B;K)$, is
identified exactly from the augmented benchmark, for every admissible $K$.
\item \emph{First-order failure certificate.} If $D_B(\Delta;K_0)\neq0$ for
some $\Delta\in\cD(A_{\rm aug},h)$, then $\cI(B;\cdot)$ is not identified near
$K_0$: for every sufficiently small $\varepsilon>0$ the correlation matrices
$K_0\pm\varepsilon\Delta$ induce the same law of $(Y_{A_{\rm aug}},\Theta)$ and
different values of $B$.
\end{enumerate}
\end{proposition}

The proof, including the row-space factorisation used in part (i), is given in
Appendix~\ref{sec:app-identification}.

By contraposition, local identification requires $D_B(\cdot;K_0)$ to vanish on
$\cD(A_{\rm aug},h)$. This condition need not be sufficient, just as the
row-space condition is sufficient but need not be necessary. Both checks are
linear algebra: compute a row-space inclusion and the kernel in
\eqref{eq:joint-invisible}, then restrict the linear functional $D_B$ to that
kernel.

The number of additional measurements required depends on the covariance
restrictions. In the four-point
stationary model of \cref{thm:minimal}, augmenting $\{Z_0\}$ with $Z_1$ reduces
the stationary invisible space from dimension $\numAugStatBefore$ to
$\numAugStatAfter$ even though the span condition for $B=\{Z_1,Z_2\}$ fails:
identification there comes from stationarity, not from
\cref{prop:augmentation}(i). Without that restriction the invisible dimensions
for $\{Z_0\}$, $\{Z_0,Z_1\}$ and $\{Z_0,Z_1,Z_2\}$ are $\numAugFreeBefore$,
$\numAugFreeOne$ and $\numAugFreeTwo$: observing $Z_1$ leaves a two-dimensional
invisible space, whereas observing both $Z_1$ and $Z_2$ removes all invisible
directions. Under stationarity, observing $Z_1$ already suffices. For an affine
structured covariance class with direction space $\cT$, the same local test
applies on $\cD(A_{\rm aug},h)\cap\cT$; stationarity is the Toeplitz instance.

A calibration protocol of full column rank satisfies the span condition and
removes every invisible direction in the unrestricted model. Full column rank
is stronger than is needed to identify one specified alternative, but it gives
a simple route to evaluating an entire candidate family. The next question is
how accurately that route works when dense calibration is available for only
$m$ units.
